% from main.tex
%%%%%%%%%%  scan idx 378  |  folio 365  %%%%%%%%%%
\section*{Exposé VI. Le \(K^{\bullet}\) d'un fibré projectif : calculs et conséquences}
% Exposé VI: the print's doubly underlined S and underlined Proj (scan check, 10 Nov)
\renewcommand{\Sym}{\mathbf{S}}
\renewcommand{\Projz}{\underline{\operatorname{Proj}}_0}
\begin{center}
par Pierre Berthelot
\end{center}

\section*{1. Calcul du \(K^{\bullet}\) d'un fibré projectif : cas des faisceaux localement libres de type fini}

Dans ce paragraphe, la notation \(K^{\bullet}(X)\) désignera le groupe de Grothendieck des faisceaux localement libres de type fini sur un schéma \(X\).

Le but du paragraphe est de démontrer le théorème suivant :

\textbf{Théorème 1.1.} Soient \(S\) un schéma quasi-compact, \(E\) un \(\mathcal O_S\)-Module localement libre de rang \(r+1\), \(X=\mathbf P(E)\) le fibré projectif associé ; le morphisme structural : \(X\xrightarrow{f}S\) définit un homomorphisme d'anneaux :
\[
  K^{\bullet}(S) \xrightarrow{f^*} K^{\bullet}(X)
\]
munissant \(K^{\bullet}(X)\) d'une structure de \(K^{\bullet}(S)\)-algèbre. Soit \(\xi\in K^{\bullet}(X)\) la classe du faisceau \(\mathcal O_X(1)\). Alors \(K^{\bullet}(X)\) est un \(K^{\bullet}(S)\)-module libre admettant pour base \(1,\ldots,\xi^r\).

\textbf{Plan de la démonstration :}

\(1^\circ\)) Les puissances de \(\xi\) engendrent \(K^{\bullet}(X)\) comme \(K^{\bullet}(S)\)-module.

Pour le montrer, nous allons montrer les points suivants :

\begin{enumerate}[label=\alph*)]
\item Si \(F\) est un \(\mathcal O_X\)-Module localement libre de type fini, il existe un \(\Sym(E)\)-Module module gradué \(M\) de présentation finie, plat sur \(\mathcal O_S\), tel que \(F=\Projz(M)\) (\(\Sym(E)\) étant l'algèbre symétrique de \(E\)~; pour la notation \(\Projz\), voir EGA II 8.12.1.

\item Il existe une résolution graduée finie de \(M\) par des \(\Sym(E)\)-Modules gradués \(L_i\) localement libres de type fini : \(0\longrightarrow L_{r+1}\longrightarrow\cdots\longrightarrow L_0\longrightarrow M\longrightarrow 0\).

\item Les \(L_i\) admettent une filtration finie dont les quotients sont du type \(N_j\otimes_{\mathcal O_S}\Sym(E)(j)\), les \(N_j\) étant localement libres de type fini sur \(S\).
\end{enumerate}

